\section{Exact optimization of a bounded crossing window}
\label{sec:ordering}

The previous scan-order planner is a fast greedy search with several starting
crossings and tie rules. We add a deterministic local optimization with a
precise objective. It is independent of the algebraic quotient and remains
optional, since a better frontier profile need not yield a faster homology
calculation.

\subsection{The subset state is exact}
Let $P$ be the fixed set of crossings before a window, and let $B$ be the
$k$ crossings inside it.
For $S\subseteq B$, let $\partial(P\cup S)$ be the set of edge labels occurring
an odd number of times among crossings already processed. Write
\[
 b(S)=|\partial(P\cup S)|.
\]
Every edge label occurs exactly twice in the full diagram. Thus this set is
exactly the scan frontier, including the cancellation of labels that occur
twice at a self-loop crossing. It depends on the subset $S$, not the order
used to reach $S$.

Only the at most $4k$ labels incident with the window can change. Represent
those labels by bit positions, let $z$ encode the initial local frontier, and
let $t_i$ be the parity mask of crossing $i$. If $c$ counts frontier labels
outside this local set, then
\begin{equation}\label{eq:subsetwidth}
 b(S)=c+\operatorname{popcount}\left(z\mathbin{\mathrm{xor}}
                    \mathop{\mathrm{xor}}_{i\in S}t_i\right).
\end{equation}
All subset widths can be computed using the least significant set bit of
each subset. The sets at both ends of the window are fixed, so every frontier
outside the window is unchanged by its permutation.

\subsection{Why the optimization has two passes}
The selected objective is lexicographic: first minimize the global maximum
frontier size, then an additive cost. The implementation offers either
$c_1(b)=b$ or $c_2(b)=2^{b/2}$. Frontier sizes are even for a degree-four
diagram graph. The latter cost is motivated by the dimension of a dot
endomorphism space, and is explicitly only a proxy for full scanner work.

A tempting one-label recurrence for pairs $(\text{peak},\text{sum})$ is
incorrect. A path with a smaller early peak but a larger sum can lose to a
different path when a later frontier raises both peaks to the same value.
Lexicographic comparison is not preserved by the operation that takes a
maximum in the first coordinate. Our implementation first solves the
bottleneck problem and then the constrained additive problem.

Let $W_{\mathrm{out}}$ be the maximum frontier size outside the window. Set
$p(\varnothing)=W_{\mathrm{out}}$ and
\[
 p(S)=\min_{i\in S}\max\{p(S\setminus\{i\}),b(S)\}.
\]
The optimal global peak is $W_*=p(B)$. Now set $d(\varnothing)=0$, forbid
every nonempty state with $b(S)>W_*$, and compute
\[
 d(S)=c_j(b(S))+\min_{i\in S}d(S\setminus\{i\})
\]
over feasible predecessors. Backtracking gives an optimal window order.
The original order is feasible, and equal-score replacements are discarded
to avoid unnecessary changes.

\begin{theorem}[Exact window optimization]\label{thm:window}
For fixed crossings outside a window of size $k$, the two-pass algorithm
finds a permutation minimizing the stated global lexicographic objective.
It uses $O(k2^k)$ subset transitions and $O(2^k)$ table entries, in addition
to polynomial work for boundary masks and exact integer costs. Every
accepted replacement weakly decreases the global peak and strictly improves
the selected global pair.
\end{theorem}
\begin{proof}
Equation~\eqref{eq:subsetwidth} gives the exact cost at every state. Every
window permutation corresponds to exactly one maximal path through the
subset lattice. Conversely, each such path adds one distinct crossing at
each step and yields a permutation. The first recurrence is the usual
bottleneck recurrence on this acyclic graph, with all fixed outside peaks
included in its initial value. After $W_*$ is known, the second recurrence
is an ordinary additive shortest path in the induced feasible subgraph.
Outside frontiers are fixed, so minimizing the local additive contribution
minimizes the global one. The lattice contains $2^k$ states and
$k2^{k-1}$ directed inclusion edges. Integer masks have $O(k)$ local bits;
the additive costs have at most polynomially many bits in the input size.
Thus the corresponding bit complexity is $2^k\operatorname{poly}(n,k)$.
\end{proof}

The bounded number of overlapping passes is part of the algorithm. When a
time budget expires, the planner returns the last completely verified order.
This is still a valid input to the recognizer. If $k=O(\log n)$, a polynomial
number of windows costs polynomial time; if $k=O((\log n)^2)$, their cost is
quasi-polynomial. These statements concern the planner. They do not establish
a width bound for its output, a globally optimal unrestricted order, or a
quasi-polynomial homology computation.

The implementation is independently checked against all permutations of
many small degree-four multigraphs, including loops and parallel edges.
The graph theorem does not require those test graphs to be realizable as
knot diagrams. Separate tests on validated knot diagrams verify that the
resulting orders preserve homology.
